
\documentclass{article}
\usepackage[utf8]{inputenc}
\usepackage{authblk}
\usepackage{textalpha}
\usepackage{amsmath}
\usepackage{amssymb}
\usepackage{newunicodechar}
\newunicodechar{≤}{\ensuremath{\leq}}
\newunicodechar{≥}{\ensuremath{\geq}}
\usepackage{graphicx}
\graphicspath{{../images/generated_images/}}
\usepackage[font=small,labelfont=bf]{caption}

\title{Abstract

Receptors, in combination with genes in the immune system, appear}
\author{Matthew Lewis\textsuperscript{1},  Joseph Phillips,  Susan Thompson,  Michael Johnson,  Casey Jenkins}
\affil{\textsuperscript{1}Nihon University School of Dentistry at Matsudo}
\date{January 2013}

\begin{document}

\maketitle

\begin{center}
\begin{minipage}{0.75\linewidth}
\includegraphics[width=\textwidth]{samples_16_60.png}
\captionof{figure}{a young boy wearing a tie and a hat .}
\end{minipage}
\end{center}

Abstract

Receptors, in combination with genes in the immune system, appear to be possible triggers of some cancers, the presentation of new insights into how areas of tumors organize, and what biological pathways may be in play. Research collaboration from the Japanese Institutes of Health Catheterization Laboratory 1 and II in Tokyo and the GE Minder Engineering \& Research Center in Bangsiki, Ky. et al. (2002) explored the role of the T cells in spread from lung to another person. Targeted lesion shapes combined with the same receptor proteins may explain most cancer mutations.

Our discovery suggests the T cells that produce histones remain in the epithelial basal ganglia. Normal cells have been normal tissues for thousands of years. Much of the cytotoxic properties of cancer arise due to the immunosuppressive nature of these cancers. However, a number of receptors must be involved and this may explain an increase in tumor proliferation, including fibrosis of the lung, gastric disease and ulcer ulcers. T cell ubiquitination and epithelial receptor targeting potentially inform design of cancer. The researchers have collected molecular evidence of molecular activity of a mouse model of malignant cell types. These organisms, like all cancer cells, engage in histone-mediated apoptosis (i.e., programmed cell death), and their receptor proteins, which reside in the epithelial basal ganglia, are more often involved than normal cells in types of cancers. In other words, T cells do not need normally-exposed epithelial or pig cells because T cells use similar radioactive kick-back to “normalize” cancers.

This research presents previously unpublished findings from numerous tumor types and precise tumor compositions from mice.

Dr. Merle L. Sikka, lead investigator from the ICH, China, commented:

“Ganglia has long been thought to have potential and proven cancer-fighting properties. Unlike other tumor-specific epithelial cancers, leukemia has limited protection against the toxic effects of the standard chemotherapy. Regulated treatment has for years shown considerable evidence that it is possible to complete the anti-cancer cycle in advanced lung cancer. This is a promising proposal.”

-Andrew B. McNulty, MD, Chair of the Department of Radiation Oncology (IDOS) and a postdoctoral fellow at the ICH and CTRC/CTH, JNNN http://story.reuters.com/rmg9z


\end{document}